\begin{center}
    {\LARGE \textbf{ĐỀ LUYỆN 05}}\\[6pt]
    {\large (Dành cho đội tuyển quốc gia)}
\end{center}

\begin{tabularx}{\textwidth}{|c|X|c|c|c|}
\hline
\textbf{Bài} & \textbf{Tên bài} & \textbf{File đầu vào} & \textbf{File đầu ra} & \textbf{Điểm} \\
\hline
A & Tổng số trong xâu & \texttt{A.INP} & \texttt{A.OUT} & 6 \\
\hline
B & Số hoàn hảo kiểu mới & \texttt{B.INP} & \texttt{B.OUT} & 5 \\
\hline
C & Tránh đối xứng K & \texttt{C.INP} & \texttt{C.OUT} & 5 \\
\hline
D & Ước bội chính phương & \texttt{D.INP} & \texttt{D.OUT} & 4 \\
\hline
\end{tabularx}

\subsection*{Bài A}

Cho một xâu ký tự bao gồm các chữ cái latin và chữ số, có độ dài không quá $10^6$.

Yêu cầu: Tính tổng các số trong xâu đã cho (mỗi số trong xâu có giá trị không quá $10^6$).

\begin{itemize}
    \item Nếu trong xâu không có ký tự chữ số nào thì ghi $-1$.
\end{itemize}

\textbf{Dữ liệu vào}
\begin{itemize}
    \item Một dòng duy nhất chứa xâu ký tự.
\end{itemize}

\textbf{Dữ liệu ra}
\begin{itemize}
    \item Một số là tổng các số trong xâu hoặc $-1$.
\end{itemize}

\textbf{Ví dụ}
\begin{itemize}
    \item Input: \texttt{ktab3b6afh12} $\rightarrow$ Output: $21$
    \item Input: \texttt{qutanyh} $\rightarrow$ Output: $-1$
\end{itemize}

\vspace{0.5cm}

\subsection*{Bài B}

Một số được gọi là \textbf{số hoàn hảo kiểu mới} nếu nó không chia hết cho số nguyên tố nào lớn hơn $10$.

Cho số nguyên $n \leq 10^5$ và dãy $a_1, a_2, ..., a_n$.

Hãy tính tổng các số hoàn hảo kiểu mới.

\textbf{Dữ liệu vào}
\begin{itemize}
    \item Dòng 1: $n$
    \item Dòng 2: $a_1, a_2, ..., a_n$
\end{itemize}

\textbf{Dữ liệu ra}
\begin{itemize}
    \item Tổng các số thỏa mãn.
\end{itemize}

\textbf{Ràng buộc}
\begin{itemize}
    \item $1 \le a_i \le 10^7$ (30\%)
    \item $1 \le a_i \le 10^{18}$ (70\%)
\end{itemize}

\textbf{Ví dụ}
\begin{itemize}
    \item Input: $5$, $9\ 100\ 10\ 26\ 30$ $\rightarrow$ Output: $149$
    \item Input: $2$, $100000000\ 81$ $\rightarrow$ Output: $100000081$
\end{itemize}

\vspace{0.5cm}

\subsection*{Bài C}

Cho chuỗi $S$ độ dài $N$ chỉ gồm chữ cái thường.

Hãy đếm số chuỗi hoán vị của $S$ sao cho \textbf{không chứa} chuỗi con đối xứng độ dài $K$.

\textbf{Dữ liệu vào}
\begin{itemize}
    \item Dòng 1: $N, K$
    \item Dòng 2: $S$
\end{itemize}

\textbf{Dữ liệu ra}
\begin{itemize}
    \item Số lượng chuỗi hợp lệ.
\end{itemize}

\textbf{Giới hạn}
\begin{itemize}
    \item $2 \le K \le N \le 10$
\end{itemize}

\textbf{Ví dụ}
\begin{itemize}
    \item Input: 
    \begin{verbatim}
    3 2
    aab
    \end{verbatim}
    Output: $1$

    \item Input:
    \begin{verbatim}
    5 3
    zzyyx
    \end{verbatim}
    Output: $16$
\end{itemize}

\vspace{0.5cm}

\subsection*{Bài D}

Cho dãy số nguyên dương $A$ gồm $N$ phần tử.

Đếm số cặp $(i, j)$ với $1 \le i < j \le N$ sao cho:

$$
\frac{\mathrm{BCNN}(a_i, a_j)}{\mathrm{UCLN}(a_i, a_j)}
$$

là một số chính phương.

\textbf{Dữ liệu vào}
\begin{itemize}
    \item Dòng 1: $N$
    \item Dòng 2: $a_1, a_2, ..., a_N$
\end{itemize}

\textbf{Dữ liệu ra}
\begin{itemize}
    \item Số cặp thỏa mãn.
\end{itemize}

\textbf{Giới hạn}
\begin{itemize}
    \item $1 \le N \le 10^6$
    \item $1 \le a_i \le 10^6$
\end{itemize}

\textbf{Ví dụ}
\begin{verbatim}
Input:
5
2 8 3 12 50

Output:
4
\end{verbatim}

\textbf{Giải thích}

Các cặp hợp lệ là: $(2,8), (2,50), (8,50), (3,12)$.